\input{../tex-inputs/latex-leseansicht-vorspann}

\section[Arthur Schnitzler an Gustav Schwarzkopf, {[}zwischen 1893 und Mitte 1894?{]}]{L03971 Arthur Schnitzler an Gustav Schwarzkopf, {[}zwischen 1893 und Mitte 1894?{]}}
\nopagebreak\mylabel{L03971v}
\rehead{ }\normalsize\beginnumbering\briefempfaengerindex{Schwarzkopf, Gustav@\textsc{Schwarzkopf, Gustav}!zzzSchnitzler, Arthur@\emph{von Arthur Schnitzler}!1894-06-301@{{[}zwischen 1893 und Mitte 1894?{]}}|(be}
\toendnotes[C]{\smallbreak\pagebreak[2]}
\correspDesc{Versand  durch Arthur Schnitzler im Zeitraum [zwischen 1893 und  Mitte 1894?] in Wien
\newline{}Erhalt  durch Gustav Schwarzkopf im Zeitraum [zwischen 1893 und  Mitte 1894?] in Wien}\toendnotes[C]{\smallbreak}
\Standort{CUL, Schnitzler, B 96.}
\physDesc{Brief, 1 Blatt, 1 Seite, 156 Zeichen
\newline{}Handschrift: Bleistift, deutsche Kurrent}\toendnotes[C]{\smallbreak}
\pstart
           {\pb}\uline{Do{\geminationn}erſtag}\hfill Mittg.\pend
           
\pstart{}\label{K_L03971-1v}\edtext{Verehrteſter Freund}{\lemma{\textnormal{\emph{Verehrtester Freund}}}\Cendnote{\textnormal{Das Blatt ist undatiert. Die
               verwendete Grußformel deutet aber auf den Zeitraum 1893 bis Mitte 1894 hin, in dem 
               Schnitzler häufig Anreden an Schwarzkopf\pwindex{Schwarzkopf, Gustav 7.\,11.\,1853 Wien – 13.\,11.\,1939 ebd.@\textsc{Schwarzkopf, Gustav} (7.\,11.\,1853 Wien – 13.\,11.\,1939 ebd.)|pwk} mit »Verehrter« und
               »Verehrtester« verwendet.}}}\label{K_L03971-1},\pend\vspace{0.5em}
\pstart
           leider kann ich Sie heute nicht abholen! Auf hoffentlich ſehr baldiges
               Wiederſehen,\pend
           
\pstart
           Ihr Sie hochſchätzender{\\[\baselineskip]}\spacefill\mbox{Arth Schnitzler}\pend
           \leftskip=0em{}\selectlanguage{ngerman}\endnumbering\briefempfaengerindex{Schwarzkopf, Gustav@\textsc{Schwarzkopf, Gustav}!zzzSchnitzler, Arthur@\emph{von Arthur Schnitzler}!1893-01-011@{{[}zwischen 1893 und Mitte 1894?{]}}|)be}\mylabel{L03971h}
\begin{anhang}
\end{anhang}\newcommand{\dateiname}{L03971}\newcommand{\titel}{Arthur Schnitzler an Gustav Schwarzkopf, [zwischen 1893 und Mitte 1894?]}\newcommand{\editorInnen}{Herausgegeben von Jahnke, SelmaMüller, Martin Anton}\input{../tex-inputs/latex-leseansicht-abspann}
